\section{引言}
\label{sec:intro}
共享微电网上的三栋建筑各自运行一个能源控制器，在七个目标（成本、碳排放、峰值、不舒适度、设备损耗、爬坡和韧性）之间权衡。
我们调节这 $3\times7$ 个目标权重，以最小化区域净负荷的爬坡量，它由一个确定性的黑箱仿真器测得。评分分为两个阶段：\emph{离线}
阶段有 2,048 行实测数据和两次提交机会；\emph{在线}阶段，一个冻结的程序最多可以顺序调用仿真器 20 次，每次最多 50 个点，且在
首次调用前只能使用训练行。

在线阶段有三个特点，使标准的贝叶斯优化~\cite{rasmussen,turbo}难以奏效。搜索空间是 21 维、受单纯形约束的，预算只有 1,000 个
点，而用来给候选排序的代理模型恰恰在改进所在之处不可靠：在离线阶段我们能测到的两个点上，代理模型的预测偏乐观 8--43~kWh
（第~\ref{sec:problem}~节）。另一方面，每次调用都返回 50 个精确测量值，且以最小值计分，因此可以把一个批次设计成一组实验，
而不是一串预测的最优解。我们的贡献正是利用了这一点：
\begin{itemize}
\item \textbf{加性分块下降}（第~\ref{sec:block}~节）。每个候选只改动一栋建筑的权重，因此其效果被单独测得；下一次调用把各建筑
的改进组合起来，组合的值由加性模型仅凭实测效果预测，无需拟合模型。
\item \textbf{受控的评估}（第~\ref{sec:exp}~节）。在课程发布的仿真器上做完整的 20 次调用预演，每个主要变体六个种子，用精确
置换检验，并设计一个只去掉分块结构的消融。我们也报告一个负结果：证据驱动的名额分配（第~\ref{sec:alloc}~节）按实测命中率
重新分配批次名额，它改变了点的去向，却没有改变最终 $R$。
\end{itemize}
策略被部署为一个确定性的运行包，使被评分的运行等于其预演。
